\chapter{Introduction}

``Digital security is the computer equivalent of disaster insurance.
Few people care very much about it or give it much thought, and everyone hates paying for it… until a catastrophe hits.
Then we are either really glad we had it or really sad that we didn't have enough of it or didn't have it at all.''\cite{tpmbook}.

In an ideal world, there would be no malicious actors, no hackers, and cybersecurity would not need to exist.
But we live in a world where cyberattacks are a daily phenomenon, yet companies, and people in general, keep choosing convenience over security.
Now that we live in ``the age of AI'', it seems like everyone is on a hunt for personal data, from hackers to big tech companies.
We barely started to understand that the internet is not a safe place, yet we find ourselves in a world where attackers are always one step ahead.
I hope that people slowly bypass their ``what do I have to hide?'' era, and start taking measures for their digital security.

There are clear measures, but they are usually inconvenient.
It is hard to make something hard for an attacker, but frictionless for the user.
Passwords are still the most common authentication and authorization method, although they have been proven inefficient since their birth.
We can see more and more companies enforcing two-factor authentication (2FA) on the users of their applications.
This way, the password gets paired with, for example, something the user has.
Now the black hat has to steal your phone as well, not just guess your dog's name and year of birth.

Digital good practices involve choosing security measures like 2FA, but not only.
One ground rule should be: never use company devices for personal use, and vice versa.
If each employee used their work laptop for work only, probably the phishing statistics would look better as well.
And yet, things get mixed very early on: you start to work at company Y (to avoid the saying X), and the first thing is that you get an account which needs a mobile authenticator app registered.
You pull out your phone, scan the QR code, and now you have your personal device coupled with your work life.

This coupling is a light one, but my thesis tries to break it.
I will present you an alternative to regular authenticator apps, which does not use your personal device.
And as long as you follow my guidelines, it won't be less secure.
